\section{Preliminaries}
\paragraph{\textit{Definition 1. (Cell Complex~\cite{HGh19})}} 
A \textbf{regular cell complex} is a topological space $X$ decomposed into a collection of disjoint subspaces $\{x_\alpha\}_{\alpha \in P_X}$, referred to as \textit{cells}, satisfying the following conditions:
% \vspace{-5pt}
\begin{enumerate}[left=1.2em, label=\arabic*.]
    \item For each point $p \in X$, there exists an open neighborhood intersecting only finitely many cells.
    \item For any pair of cells $x_\alpha, x_\tau$, the intersection $x_\tau \cap \overline{x_\alpha}$ is nonempty if and only if $x_\tau \subseteq \overline{x_\alpha}$, where $\overline{x_\alpha}$ denotes the topological closure of $x_\alpha$.
    \item Each cell $x_\alpha$ is homeomorphic to an open ball in $\mathbb{R}^n$ for some non-negative integer $n$.
    \item \textit{(Regularity)} The closure $\overline{x_\alpha}$ of every cell is homeomorphic to a closed ball in $\mathbb{R}^{n_\alpha}$, with the interior mapped homeomorphically onto $x_\alpha$ itself.
\end{enumerate}
% \vspace{-5pt}
\paragraph{\textit{Definition 2}} A \textbf{cellular lifting map} is a function \( f: \mathcal{G} \to X \) from the space of graphs \(\mathcal{G}\) to the space of regular cell complexes \(X\), satisfying that two graphs \(G_1, G_2 \in \mathcal{G}\) are isomorphic if and only if their corresponding cell complexes \(f(G_1)\) and \(f(G_2)\) are isomorphic. Intuitively, a cell complex is built hierarchically by first considering 0-cells (vertices), then attaching 1-cells (edges) via their endpoints, and further incorporating higher-dimensional cells—such as 2-cells—by gluing disks along cycles.

% \vspace{-5pt}
\paragraph{\textit{Definition 3. (Retrieved Cellular Topologies-Augmented Graph Learning)}}  
Given a graph $\mathcal{G}$, a cellular lifting map  $f: \mathcal{G} \to X$ is conducted from the space of graphs to the space of regular cell complexes. 
To construct a cell complex knowledge base \(X = \{X_\sigma\}\), we extract localized subcomplexes from the evolving complex \(X\) centered around a master 0-cell \( x_m^0 \), defined as the 0-cell with the highest degree in the complex. Each subcomplex includes the set \(\mathrm{k\text{-}hop}_{X^{(0)} \cup X^{(1)}}(x_m^0)\) of 0- and 1-cells reachable within \(k\) hops from \(x_m^0\), as well as all 2-cells \(x^2 \in X^{(2)}\) such that there exists a cell \(x^1 \in \mathrm{k\text{-}hop}_{X^{(0)} \cup X^{(1)}}(x_m^0)\) with \(x^1 \prec x^2\), meaning \(x^1\) is a face of \(x^2\). To enable retrieval, each cell complex $X_\sigma$ is indexed by a composite \textbf{key} that includes: the timestamp $\tau_b$, the master 0-cell embedding $h_{m}^0$,  and a two-dimensional topological characteristic $h_{\sigma}^2$ obtained by pooling over all 2-cells within $X_\sigma$. Retrieval is performed via similarity over keys to obtain \textbf{values} such as task-specific output vectors $\{o_i \in \mathbb{R} \mid x_i \in X_\sigma\}$ and cell complex embeddings $\{h_i \in \mathbb{R} \mid x_i \in  X_\sigma\}$. 

Given a  graph $\mathcal{G}$, we first split it into training and testing subsets, $\mathcal{G} = \mathcal{G}_{\rm{train}} \cup \mathcal{G}_{\rm{test}}$, and lift them into a regular cell complex $X = X_{\rm{train}} \cup X_{\rm{test}}$ via a map $f: \mathcal{G} \rightarrow X$. Each unit (e.g., a 0-cell $x^0_i$, 1-cell $x^1_{ij}$, or complex $X_\sigma$) has label $y_i$ if and only if it resides in $X_{\rm{train}}$. The goal is to predict $Y_{\rm{test}}$ by retrieving relevant cell complexes and propagating structure-aware representations. %Following RAGraph~\cite{JiangQXZZ00X0W24}, we unify node-, edge-, and graph-level tasks within a similarity-based framework. Each instance produces a task-specific output vector \( o \), whose label is predicted by comparing it with the ground truth or a class prototype. In the \(\kappa\)-shot node/graph classification setting, each class \( c \in \mathcal{C} \) is associated with a prototype vector computed from the support set \( \mathcal{D} \) as \( \tilde{o}_c = \frac{1}{\kappa} \sum_{(i, y_i) \in \mathcal{D},\, y_i = c} o_i \), and the prediction is made by \( y_i = \arg\max_{c \in \mathcal{C}} \texttt{sim}(o_i, \tilde{o}_c) \).
%For edge prediction, a link between \( v_i \) and \( v_q \) is inferred if \( \texttt{sim}(o_i, o_q) \geq \texttt{sim}(o_i, o_j) + \epsilon \) for any \( (v_i, v_j) \in \mathcal{E}_{\text{train}} \). 